\documentclass[11pt]{article}
\usepackage[a4paper,margin=26mm]{geometry}
\usepackage[T1]{fontenc}
\usepackage{lmodern}
\usepackage{amsmath,amssymb,amsthm}
\usepackage{microtype}
\usepackage{xurl}
\usepackage{hyperref}
\hypersetup{colorlinks=true,linkcolor=blue,urlcolor=blue,citecolor=blue,
 pdftitle={The decomposition lengths of a numerical semigroup form an interval},
 pdfauthor={Horace Dong}}
\newtheorem{theorem}{Theorem}
\newtheorem{lemma}{Lemma}
\newtheorem{proposition}{Proposition}
\newtheorem{corollary}{Corollary}
\theoremstyle{definition}
\newtheorem{definition}{Definition}
\theoremstyle{remark}
\newtheorem{remark}{Remark}
\newcommand{\repoversion}{1.24.0}
\setlength{\emergencystretch}{1.5em}
\newcommand{\N}{\mathbb N}
\DeclareMathOperator{\SG}{SG}
\DeclareMathOperator{\G}{G}
\newcommand{\F}{\mathrm F}

\title{The decomposition lengths of a numerical semigroup\\ form an interval}
\author{Horace Dong\\ \small Independent researcher\\
\small\texttt{maxwelldhx@gmail.com}\\
\small ORCID: \href{https://orcid.org/0009-0002-7554-7548}{0009-0002-7554-7548}}
\date{20 November 2026}
\begin{document}
\maketitle

\begin{abstract}
Every numerical semigroup is an irredundant intersection of finitely many
irreducible numerical semigroups, often in several ways, and the number of
components can vary. García-Sánchez conjectured that, for each numerical
semigroup, the set of these numbers is an interval, and García-Sánchez and
O'Neill proved this when the multiplicity is at most six. We prove the
conjecture for every numerical semigroup. One component can be exchanged for
another with the same Frobenius number, and when this leaves a special gap
uncovered, one component with a smaller Frobenius number repairs it. This
changes the length by at most one at each step, and every such process ends in
a normal form whose length depends only on the semigroup. We also show that a
numerical semigroup of even multiplicity $2k$ has a decomposition with at most
$k$ components, which together with the first result proves a second
conjecture of García-Sánchez and O'Neill. Both results are formally verified
in Lean~4. The proofs were found and formalized by AI agents working under the
author's direction; their role is described at the end of the paper.
\end{abstract}

\section{Introduction}

Let $\N=\{0,1,2,\dots\}$. A \emph{numerical semigroup} is a subset
$S\subseteq\N$ that contains $0$, is closed under addition and has finite
complement $\G(S)=\N\setminus S$, the set of \emph{gaps} of~$S$. If
$S\ne\N$, its \emph{Frobenius number} $\F(S)$ is the largest gap; its
\emph{multiplicity} $m(S)$ is the smallest positive element. A numerical
semigroup is \emph{irreducible} if it is not the intersection of two numerical
semigroups properly containing it.

Every numerical semigroup is the intersection of finitely many irreducible
ones~\cite[Chapter~3, Proposition~4.44]{rg}. Following
García-Sánchez~\cite{g} and García-Sánchez and O'Neill~\cite{go}, a
\emph{decomposition} of $S$ (into irreducibles) is an expression
$S=S_1\cap\dots\cap S_n$ with every $S_i$ irreducible that is
\emph{irredundant}: $S\ne\bigcap_{i\in I}S_i$ for every proper subset
$I\subsetneq\{1,\dots,n\}$. The number $n$ is the \emph{length} of the
decomposition, and we write $\mathcal L(S)$ for the set of lengths of all
decompositions of $S$.

The set $\mathcal L(S)$ is finite: a decomposition of $S$ has at most
$\#\SG(S)\le m(S)-1$ components, where $\SG(S)$ is the set of special gaps
defined in Section~\ref{sec:prelim} (see~\cite[Section~2]{go}, or
Lemma~\ref{lem:private} and the remark after Theorem~\ref{thm:even} below). García-Sánchez~\cite{g}, who calls
decompositions factorizations, showed that for every $k\ge2$ the union of all
sets $\mathcal L(S)$ containing $k$ is $\{2,3,\dots\}$; he reported that he had
found no numerical semigroup whose set of lengths is not an interval, and suggested
proving that it always is~\cite[Section~3]{g}. García-Sánchez and
O'Neill~\cite{go} state:

\begin{quote}
\textbf{Conjecture 1.1 of~\cite{go}.} The set of decomposition lengths of any
numerical semigroup is an interval.
\end{quote}

\noindent They proved it for multiplicity at most six~\cite[Section~3]{go},
found large intervals of lengths for ordinary numerical
semigroups~\cite[Section~4]{go}, and proposed a refinement for even
multiplicity:

\begin{quote}
\textbf{Conjecture 3.9 of~\cite{go}.} Fix a numerical semigroup $S$ of
multiplicity $m=2k$ with $k\in\mathbb Z$. If $S$ has a decomposition of
length $j$ with $k<j<m$, then $S$ has a decomposition of length $j-1$.
\end{quote}

\noindent For $m=6$ this is their Theorem~3.5, and they report computational
evidence for $m=8$. We prove both conjectures.

For $S\ne\N$ let
\[
r(S)=\#\{g\in\SG(S): 2g>\F(S)\}
\]
be the number of special gaps in the upper half. (The empty family has
intersection $\N$ by convention, so $\mathcal L(\N)=\{0\}$; with the
alternative convention that $\N$, being irreducible, is its own decomposition
of length $1$, one gets $\{1\}$. Either way it is an interval; for $S\ne\N$
the two conventions agree.)

\begin{theorem}\label{thm:interval}
For every numerical semigroup $S$, the set $\mathcal L(S)$ is a nonempty set
of consecutive integers: if $a,b\in\mathcal L(S)$ and $a\le j\le b$, then
$j\in\mathcal L(S)$.
\end{theorem}

\begin{theorem}\label{thm:even}
Let $S$ be a numerical semigroup of multiplicity $2k$. Then $S$ has a
decomposition of length at most $k$.
\end{theorem}

\begin{corollary}\label{cor:descent}
Conjecture~3.9 of~\cite{go} holds. More precisely, if $S$ has multiplicity
$2k$ and a decomposition of length $j>k$, then it has a decomposition of
length $j-1$.
\end{corollary}

\begin{proof}
By Theorem~\ref{thm:even}, $\min\mathcal L(S)\le k\le j-1<j$, so $j-1$ lies
between two elements of $\mathcal L(S)$ and belongs to $\mathcal L(S)$ by
Theorem~\ref{thm:interval}.
\end{proof}

The proof of Theorem~\ref{thm:interval} is elementary. It uses the
description of decompositions by special gaps from the book of Rosales and
García-Sánchez~\cite{rg} and an exchange that keeps the Frobenius number of
an irreducible semigroup fixed, of the kind used by Blanco and
Rosales~\cite{br}. We apply the exchange to one component of a decomposition
at a time. When it leaves a special gap of $S$ uncovered, we add one
irreducible component with a smaller Frobenius number, and then remove
redundant components; the length changes by at most one
(Lemma~\ref{lem:step}). The process ends at a decomposition in which every
component is in a normal form, and all such decompositions have exactly $r(S)$ components (Lemma~\ref{lem:normal});
in particular $r(S)\in\mathcal L(S)$ for $S\ne\N$ (Remark~\ref{rem:r}). Starting from a shortest and from
a longest decomposition then gives every length in between. Theorem~\ref{thm:even}
generalizes the pairing of residue classes in the proof
of~\cite[Theorem~3.5]{go}.

\subsection*{Prior work}
The existence of decompositions, the criterion by special gaps
(Lemma~\ref{lem:cover}) and an algorithm that finds a decomposition of
minimum length are in~\cite[Chapter~3, Section~4]{rg}. Rosales and
Branco~\cite{rb} gave bounds on the minimum length; as quoted by Bogart and
Seyed Fakhari~\cite[Theorem~2.1]{bsf}, from Corollaries~10 and~18
of~\cite{rb}, every numerical semigroup has a
decomposition with at most $r(S)$ components (the pseudo-Frobenius numbers
larger than $\F(S)/2$ used there are exactly the special gaps counted by
$r(S)$). We have not seen the full text of~\cite{rb}, which may also contain a
decomposition of length exactly $r(S)$; we do not claim this part of
Lemma~\ref{lem:normal} as new. Bogart and Seyed Fakhari~\cite{bsf} showed that
the minimum length is unbounded. In~\cite[Theorem~4.3]{go} decompositions of
the ordinary semigroups $\{0,m,m+1,\dots\}$ are changed step by step, each step
keeping the number of components or adding one, towards a decomposition
indexed by the special gaps; our argument works with the same kind of steps
for arbitrary $S$. The exchange of Lemma~\ref{lem:swap} is a special case of
the operations of Blanco and Rosales~\cite[Proposition~5]{br}, who build the
tree of irreducible numerical semigroups with a fixed Frobenius number. As far
as we know, the general cases of the two conjectures have not been treated
elsewhere; the searches are described in the final section.

\section{Special gaps and covers}\label{sec:prelim}

Throughout, $S\ne\N$ is a numerical semigroup. An \emph{oversemigroup} of $S$
is a numerical semigroup $T\supseteq S$; there are only finitely many, since
$T\setminus S\subseteq\G(S)$. A gap $g$ of $S$ is \emph{special} if
$S\cup\{g\}$ is a numerical semigroup, that is, if $2g\in S$ and $g+s\in S$
for every $s\in S\setminus\{0\}$; we write $\SG(S)$ for the set of special
gaps. The Frobenius number is special.

\begin{lemma}[{\cite[Chapter~3, Lemma~4.35]{rg}}]\label{lem:max}
If $T\supsetneq S$ is an oversemigroup, then $\max(T\setminus S)\in\SG(S)$.
\end{lemma}

\begin{proof}
Let $g=\max(T\setminus S)$. For $s\in S\setminus\{0\}$ the numbers $g+s$ and
$2g$ lie in $T$ and exceed $g$, so they lie in $S$.
\end{proof}

\begin{lemma}[{\cite[Chapter~3, Proposition~4.48]{rg}}]\label{lem:cover}
Let $S_1,\dots,S_n$ be oversemigroups of $S$. Then $S=S_1\cap\dots\cap S_n$
if and only if every special gap of $S$ is a gap of some $S_i$.
\end{lemma}

\begin{proof}
A special gap of $S$ is not in $S$, hence not in the intersection if the
intersection is $S$. Conversely, if the intersection $T$ properly contains $S$,
then $\max(T\setminus S)$ is a special gap of $S$ lying in every $S_i$, by
Lemma~\ref{lem:max}.
\end{proof}

\begin{remark}\label{rem:sg}
The condition is that $\SG(S)\subseteq\bigcup_i\G(S_i)$, and it cannot be
replaced by $\SG(S)\subseteq\bigcup_i\SG(S_i)$. For example, let $S$, $U$ and
$V$ be the numerical semigroups with gaps $\{1,2,3,5\}$, $\{1,2\}$ and
$\{1,3,5\}$. Then $U$ and $V$ are irreducible and $S=U\cap V$, while
$\SG(S)=\{2,3,5\}$, $\SG(U)=\{2\}$ and $\SG(V)=\{5\}$. (The second form is
printed in~\cite[(2.5)]{go}; we use the first, as in~\cite{rg}.)
\end{remark}

For an oversemigroup $T$ of $S$ let $E(T)=\SG(S)\cap\G(T)$ be the set of
special gaps of $S$ that $T$ \emph{covers}. By Lemma~\ref{lem:cover},
oversemigroups $S_1,\dots,S_n$ intersect in $S$ exactly when their sets
$E(S_i)$ cover $\SG(S)$. A special gap $h\in E(S_i)$ is a \emph{private}
point of $S_i$ in the family if no other $S_j$ covers it.

\begin{lemma}\label{lem:private}
Let $S_1,\dots,S_n$ be oversemigroups of $S$ with intersection $S$. The
expression $S=S_1\cap\dots\cap S_n$ is irredundant if and only if every $S_i$
has a private point. In particular a decomposition of $S$ has at most
$\#\SG(S)$ components, and none of them is $\N$.
\end{lemma}

\begin{proof}
If $S_i$ has a private point $h$, then $h$ lies in every $S_j$ with $j\ne i$,
so a subfamily without $S_i$ does not intersect in $S$; every proper
subfamily omits some index, so the expression is irredundant. If $S_i$ has no
private point, the family without $S_i$ still covers $\SG(S)$ and intersects
in $S$ by Lemma~\ref{lem:cover}. Private points of different components are
different, and $E(\N)=\emptyset$.
\end{proof}

Removing components without private points one at a time, we may turn every
family of oversemigroups intersecting in $S$ into an irredundant one; a
component that has a private point keeps it when others are removed. (If the
family is a list with repetitions, repeated members have no private points
and are removed.)

We use the following classical facts about irreducible numerical semigroups
(see~\cite[Chapter~3, Theorem~4.2 and Proposition~4.4]{rg}
and~\cite[Section~1]{go}). Let $Q$ be a numerical
semigroup with Frobenius number $f$. Then $Q$ is irreducible if and only if
it is maximal among the numerical semigroups not containing $f$, if and only
if
\begin{equation}\label{eq:sym}
\text{for all integers }0\le x\le f\text{ with }2x\ne f:\qquad
x\in Q\iff f-x\notin Q .
\end{equation}
(For the first equivalence: if $Q$ is maximal and $Q=A\cap B$ with
$A,B\supsetneq Q$, then $A$ and $B$ contain $f$ by maximality, so $f\in Q$, a
contradiction. Conversely, if $Q$ is irreducible and $T\supsetneq Q$ is a
numerical semigroup with $f\notin T$, then $Q\cup\{f\}$ is a numerical
semigroup, since $f$ is the largest gap of $Q$, and $Q=T\cap(Q\cup\{f\})$ with
both properly containing $Q$, a contradiction.) Moreover, if $B$ is a numerical
semigroup with $f\notin B$, then $B\cup\{n\in\N:n>f\}$ is a numerical
semigroup not containing $f$, and among the finitely many numerical semigroups
containing it and not containing $f$ there is a maximal one; this is an
irreducible numerical semigroup with Frobenius number $f$ containing $B$.

\section{An exchange with fixed Frobenius number}\label{sec:swap}

\begin{definition}
Let $Q$ be an irreducible oversemigroup of $S$ with Frobenius number $f$. We
call $Q$ \emph{normal} (with respect to $S$) if every integer $x$ with
$0<x<f/2$ that lies in $Q$ also lies in $S$.
\end{definition}

\begin{lemma}\label{lem:swap}
Let $Q$ be an irreducible oversemigroup of $S$ with Frobenius number $f$ that
is not normal. Let $u$ be the smallest integer with $0<u<f/2$ and $u\in
Q\setminus S$, and let $v=f-u$. Then
$Q'=(Q\setminus\{u\})\cup\{v\}$ is an irreducible oversemigroup of $S$ with
Frobenius number $f$, and $\G(Q')=(\G(Q)\setminus\{v\})\cup\{u\}$.
\end{lemma}

\begin{proof}
First, $u$ is not a sum of two positive elements of $Q$: these would be
positive elements of $Q$ smaller than $u$, hence in $S$ by the choice of $u$,
and then $u\in S$. So $Q\setminus\{u\}$ is closed under addition. By
\eqref{eq:sym}, $v\notin Q$, and $v>f/2>u$. We check that sums involving $v$
lie in $Q'$. We have $2v>f$. Let $q\in Q\setminus\{u\}$ be positive. If
$q>u$, then $v+q>f$. If $q<u$, then $u-q\notin Q$ (otherwise
$u=q+(u-q)$), and $0<u-q<f/2$, so $v+q=f-(u-q)\in Q$ by~\eqref{eq:sym}; as
$v+q>u$, it lies in $Q'$. Hence $Q'$ is a numerical semigroup. It contains
$S$ because $u\notin S$, its largest gap is still $f$, and it
satisfies~\eqref{eq:sym} because only the pair $\{u,f-u\}$ has changed. So
$Q'$ is irreducible.
\end{proof}

Since $u\in\G(Q')$, the exchange lowers the number of integers below $f/2$
that lie in $Q\setminus S$ by one; repeating it with the same $f$ ends at a
normal semigroup. For a decomposition of $S$ the exchange may destroy the
cover of $\SG(S)$, and we repair it as follows.

\section{Exchange and repair in a decomposition}\label{sec:step}

\begin{lemma}\label{lem:step}
Let $D$ be a decomposition of $S$ with $n$ components, one of which, $Q$, is
not normal. Then $S$ has a decomposition $D'$ with $n-1$, $n$ or $n+1$
components, obtained from $D$ by replacing $Q$ by the semigroup $Q'$ of
Lemma~\ref{lem:swap}, possibly adding one irreducible oversemigroup $R$ of $S$
with $R\supsetneq Q$, and removing components without private points.
\end{lemma}

\begin{proof}
Let $f$, $u$ and $v$ be as in Lemma~\ref{lem:swap}. By the lemma,
$E(Q')=(E(Q)\setminus\{v\})\cup(\{u\}\cap\SG(S))$: compared with $Q$, the new
component may cover $u$ and no longer covers $v$.

\emph{Case 1: $v$ is not a private point of $Q$ in $D$.} Then the family
$D_1$ obtained by replacing $Q$ by $Q'$ still covers $\SG(S)$, so it
intersects in $S$. The private points of $Q$ in $D$ are different from $v$
and are still covered by $Q'$ and by no other component, so $Q'$ has a private
point. Another component $P$ can lose private points only if $u$ was one of
them; this happens for at most one $P$, since $u$ was private to at most one
component. Removing components without private points therefore removes at
most one component, and gives a decomposition with $n$ or $n-1$ components.

\emph{Case 2: $v$ is a private point of $Q$ in $D$.} Then $v\in\SG(S)$ is not
covered by $D_1$. Let $B=Q\cup\{x\in\N:x>v\}$, a numerical semigroup not
containing $v$, and let $R\supseteq B$ be an irreducible numerical semigroup
with Frobenius number $v$ (Section~\ref{sec:prelim}). Since $v<f$ and
$f\notin Q$, we have $R\supsetneq Q\supseteq S$, so $E(R)\subseteq E(Q)$, and
$v\in E(R)$. The family $D_2=D_1\cup\{R\}$ covers $\SG(S)$ and intersects in
$S$, and $v$ is a private point of $R$ in it. A component of $D_2$ other than
$Q'$ and $R$ keeps its private points from $D$, unless one of them is $u$,
which, as before, happens for at most one component $P$: the points covered by
$R$ were covered by $Q$, hence were not private points of other components.
So at most two components ($Q'$ and $P$) may lack private points, and removing
components without private points gives a decomposition with $n+1$, $n$ or
$n-1$ components.
\end{proof}

\begin{lemma}\label{lem:terminate}
Starting from any decomposition of $S$ and applying Lemma~\ref{lem:step}
repeatedly, one reaches after finitely many steps a decomposition in which
every component is normal.
\end{lemma}

\begin{proof}
For a numerical semigroup $T$ let $b(T)=\sum_{g\in\G(T)}(g+1)$, and for a
family $D$ let $\phi(D)=\sum_{T\in D}3^{b(T)}$. In Lemma~\ref{lem:step},
$b(Q')=b(Q)-(v-u)<b(Q)$, and $b(R)<b(Q)$ because $R\supsetneq Q$. So
replacing $Q$ by $Q'$ and $R$ replaces the term $3^{b(Q)}$ by at most
$2\cdot3^{b(Q)-1}<3^{b(Q)}$, and removing components lowers $\phi$ further.
Hence $\phi$ decreases strictly at each step, and the process stops; it stops
only when no component is non-normal.
\end{proof}

\section{The normal form and the proof of Theorem~\ref{thm:interval}}\label{sec:normal}

\begin{lemma}\label{lem:normal}
Let $D$ be a decomposition of $S$ in which every component is normal. Then
$D$ has exactly $r(S)$ components.
\end{lemma}

\begin{proof}
Let $F=\F(S)$. First let $Q$ be a normal irreducible oversemigroup of $S$ with
Frobenius number $f$. We claim that
\[
E(Q)=\{g\in\SG(S):2g<f\}\cup(\{f\}\cap\SG(S)).
\]
Let $g\in\SG(S)$. If $2g<f$, then $g\notin Q$ since $Q$ is normal and
$g\notin S$ (and $g>0$). If $g=f$, then $g\notin Q$. If $f/2<g<f$ and
$g\notin Q$, then $f-g\in Q$ by~\eqref{eq:sym}, and $0<f-g<f/2$, so $f-g\in S$
by normality; as $g$ is special, $f=g+(f-g)\in S\subseteq Q$, a
contradiction. The case $2g=f$ does not occur, since $2g\in S\subseteq Q$ for
a special gap $g$. Finally, $g>f$ implies $g\in Q$. This proves the claim.

Since $Q\supseteq S$, every Frobenius number $f$ of a component is at most
$F$. Let $g\in\SG(S)$ with $2g>F$. Then $2g>f$ for every component, so by the
claim $g$ is covered only by components with Frobenius number $g$, and $D$
contains such a component $Q_g$. These $r(S)$ components are distinct. Since
$F$ is special and $2F>F$, the component $Q_F$ covers every special gap $g$
with $2g<F$, and no special gap satisfies $2g=F$. Hence the components $Q_g$
already cover $\SG(S)$; any further component would have no private point. So
$D=\{Q_g: g\in\SG(S),\ 2g>F\}$ has $r(S)$ components.
\end{proof}

\begin{remark}\label{rem:r}
A decomposition of $S$ exists~\cite[Chapter~3, Proposition~4.44]{rg} (after
removing redundant components), and Lemma~\ref{lem:terminate} turns it into one
with all components normal; so $r(S)\in\mathcal L(S)$ for every $S\ne\N$. This
consequence is not stated in the Lean files (Section~\ref{sec:lean}).
\end{remark}

\begin{proof}[Proof of Theorem~\ref{thm:interval}]
For $S=\N$ we have $\mathcal L(\N)=\{0\}$. Let $S\ne\N$. The set
$\mathcal L(S)$ is nonempty and finite (Lemma~\ref{lem:private}); let
$a$ and $b$ be its minimum and maximum, so $a\le r(S)\le b$ by
Remark~\ref{rem:r}. Start from a decomposition of length $a$ and apply
Lemma~\ref{lem:step} until all components are normal
(Lemma~\ref{lem:terminate}). The lengths along the way change by at most one at
each step and end at $r(S)$ (Lemma~\ref{lem:normal}), so every integer between
$a$ and $r(S)$ is the length of a decomposition. In the same way, starting from
a decomposition of length $b$ gives every integer between $r(S)$ and $b$.
\end{proof}

\section{Even multiplicity}\label{sec:even}

\begin{proof}[Proof of Theorem~\ref{thm:even}]
Let $m=2k$ be the multiplicity of $S$; then $S\ne\N$. If $g<h$ are special
gaps with $g\equiv h\pmod m$, then $h-g$ is a positive multiple of $m$, so
$h-g\in S$ and $h=g+(h-g)\in S$, a contradiction. So each residue class
modulo $m$ contains at most one special gap, and the class of $0$ contains
none. Group the residue classes into the $k$ pairs $\{i,i+k\}$,
$0\le i<k$. We show that for each pair there is an irreducible oversemigroup
of $S$ that covers all special gaps of $S$ in the pair. Then these at most
$k$ oversemigroups intersect in $S$ by Lemma~\ref{lem:cover}, and removing
redundant ones gives a decomposition with at most $k$ components.

If the pair contains no special gap, there is nothing to do. If it contains
exactly one, $f$, take an irreducible numerical semigroup with Frobenius
number $f$ containing $S$ (Section~\ref{sec:prelim}). If it contains two,
$g<f$, let $d=f-g$. Then $d\equiv k\pmod{2k}$, so $2d$ is a positive multiple
of $m$ and lies in $S$. Let $B=\{s+td:s\in S,\ t\in\N\}$, the numerical
semigroup generated by $S$ and $d$. We claim $f\notin B$. Suppose $f=s+td$.
If $t$ is even, then $td\in S$ and $f\in S$. If $t$ is odd, then
$(t-1)d\in S$ and $g=f-d=s+(t-1)d\in S$. Both are impossible. Let $R$ be an
irreducible numerical semigroup with Frobenius number $f$ containing $B$
(Section~\ref{sec:prelim}). Then $f\notin R$, and $g\notin R$ because
otherwise $f=g+d\in R$. So $R$ covers both special gaps.
\end{proof}

\begin{remark}
For $m=6$ this is the pairing of the classes $\{1,4\}$ and $\{2,5\}$ in the
proof of~\cite[Theorem~3.5]{go}. The first paragraph of the proof does not use
that $m$ is even, so $\#\SG(S)\le m-1$ for every multiplicity $m$, as noted
in~\cite[Section~2]{go}; with Lemma~\ref{lem:private}, every decomposition has
at most $m-1$ components. Theorem~\ref{thm:even} bounds only the minimum
length; the normal form of Lemma~\ref{lem:normal} can be longer than $k$. For
example, $S=\langle 4,9,10,11\rangle$ has gaps $\{1,2,3,5,6,7\}$,
$\SG(S)=\{5,6,7\}$, $r(S)=3>k=2$ and $\mathcal L(S)=\{2,3\}$.
\end{remark}

{\sloppy
\section{Formal verification}\label{sec:lean}

Theorems~\ref{thm:interval} and~\ref{thm:even} and Corollary~\ref{cor:descent}, the latter in the form of Conjecture~3.9 ($k<j<2k$), are formally verified in Lean~4 (version 4.34.1)
with Mathlib~\cite{mathlib} (commit \texttt{d13f23b7}). The $12$ files are in
version~\repoversion{} of the repository~\cite{repo}, directory
\texttt{lean/Research/TaskL/Semigroup/}.

\paragraph{The statements.}
The statement file \texttt{Challenge.lean} contains only definitions and
statements, no proofs. In ASCII notation:
{\footnotesize
\begin{verbatim}
def Numerical (S : AddSubmonoid Nat) : Prop := ((S : Set Nat)^c).Finite
def Irreducible (Q : AddSubmonoid Nat) : Prop :=
  Numerical Q /\ forall A B : AddSubmonoid Nat, Numerical A -> Numerical B ->
    Q = A inf B -> Q = A \/ Q = B
def Represents (S : AddSubmonoid Nat) (D : Finset (AddSubmonoid Nat)) : Prop :=
  forall n : Nat, n in S <-> forall Q in D, n in Q
def Decomposition (S : AddSubmonoid Nat) (D : Finset (AddSubmonoid Nat)) : Prop :=
  (forall Q in D, Irreducible Q) /\ Represents S D /\
    forall E : Finset (AddSubmonoid Nat), E ssubset D -> not Represents S E
def HasLength (S : AddSubmonoid Nat) (j : Nat) : Prop :=
  exists D : Finset (AddSubmonoid Nat), Decomposition S D /\ D.card = j
def HasMultiplicity (S : AddSubmonoid Nat) (m : Nat) : Prop :=
  0 < m /\ m in S /\ forall s in S, 0 < s -> m <= s
def LengthsExist : Prop :=
  forall S : AddSubmonoid Nat, Numerical S -> exists j, HasLength S j
def LengthsInterval : Prop :=
  forall S : AddSubmonoid Nat, Numerical S -> forall a b j : Nat,
    HasLength S a -> HasLength S b -> a <= j -> j <= b -> HasLength S j
def PositiveLengthsInterval : Prop :=
  forall S : AddSubmonoid Nat, Numerical S -> forall a b j : Nat,
    (0 < a /\ HasLength S a) -> (0 < b /\ HasLength S b) ->
    a <= j -> j <= b -> 0 < j /\ HasLength S j
def EvenMultiplicityDescent : Prop :=
  forall S : AddSubmonoid Nat, Numerical S -> forall k j : Nat,
    HasMultiplicity S (2*k) -> k < j -> j < 2*k -> HasLength S j -> HasLength S (j-1)
\end{verbatim}}
\noindent Here \texttt{\^{}c} is the complement, \texttt{inf} the
intersection of submonoids and \texttt{ssubset} strict inclusion. A numerical
semigroup is an additive submonoid of $\N$ with finite complement. Irreducibility is stated literally: if $Q=A\cap B$ for numerical
semigroups $A$ and $B$ (which then contain $Q$), then $Q=A$ or $Q=B$. A
decomposition is a finite set of irreducible numerical semigroups whose
intersection is $S$, such that no proper subset has intersection $S$; the
empty set has intersection $\N$. Using a set rather than a list loses nothing,
since a repeated component is redundant. \texttt{LengthsExist} and
\texttt{LengthsInterval} together state Conjecture~1.1 of~\cite{go} for every
numerical semigroup, including $\N$, for which the only decomposition is the
empty one; \texttt{PositiveLengthsInterval} is the same statement for positive
lengths, as in~\cite{g}; \texttt{EvenMultiplicityDescent} is Conjecture~3.9 of~\cite{go} as quoted above, with $k<j<2k$; the case $j\ge 2k$ of Corollary~\ref{cor:descent} is not stated in the Lean files and rests on the written proof above. The frozen statements do not include the bound of Theorem~\ref{thm:even},
which is proved in the Lean files on the way as \texttt{short\_\allowbreak
even\_\allowbreak decomposition}; Lemma~\ref{lem:normal} (every decomposition
with all components normal has $r(S)$ components) is proved as
\texttt{terminal\_\allowbreak card}. The assertion $r(S)\in\mathcal L(S)$ of
Remark~\ref{rem:r} and the finiteness of $\mathcal L(S)$ are not stated in the
Lean files; they rest on the written proofs.

\paragraph{Review and freezing.}
\texttt{Challenge.lean} was reviewed against~\cite{go} and~\cite{g} by a Codex
agent that had not written it, and frozen, by recording its SHA-256 hash, on
18 November 2026 at 23:47~UTC. The development logs record no compilation of a
proof file before that time, and the file was not changed afterwards.

\paragraph{The proofs.}
In \texttt{Proof.lean}, \texttt{Proof.LengthsInterval} is
\texttt{lengthsInterval} (file \texttt{Interval.lean}), which follows the
proof of Sections~\ref{sec:swap}--\ref{sec:normal}, except that the final step
is argued by contradiction: if a length $j$ between two lengths of
decompositions were missing, every decomposition would lie on the same side
of $j$ as $r(S)$ (\texttt{side\_\allowbreak of\_\allowbreak missing}, by strong
induction on $\phi$); the exchange of Lemma~\ref{lem:swap} is
\texttt{exchange\_\allowbreak of\_\allowbreak nonstandard}, the repair
component is \texttt{repair\_\allowbreak component}, the step of
Lemma~\ref{lem:step} is \texttt{rebalance} with the termination measure
$\phi$ of Lemma~\ref{lem:terminate} (\texttt{score}), and
Lemma~\ref{lem:normal} is \texttt{terminal\_\allowbreak card}.
\texttt{Proof.EvenMultiplicityDescent} is \texttt{even\_\allowbreak
of\_\allowbreak interval} applied to \texttt{lengthsInterval}, and uses
\texttt{short\_\allowbreak even\_\allowbreak decomposition}
(Theorem~\ref{thm:even}). \texttt{Proof.LengthsExist} is
\texttt{lengthsExist}, and \texttt{Proof.PositiveLengthsInterval} follows
from \texttt{lengthsInterval}. The test file \texttt{CheckSemigroup.lean}
restates the four theorems with the frozen types
(\texttt{check\_\allowbreak LengthsExist}, \texttt{check\_\allowbreak
LengthsInterval}, \texttt{check\_\allowbreak PositiveLengthsInterval},
\texttt{check\_\allowbreak EvenMultiplicityDescent}). The axioms of these
theorems, printed by \texttt{CheckSemigroup.lean}, are exactly \texttt{propext},
\texttt{Classical.choice} and \texttt{Quot.sound}. No file uses
\texttt{sorry}, \texttt{native\_\allowbreak decide},
\texttt{bv\_\allowbreak decide} or an \texttt{axiom} declaration.

\paragraph{Checks.}
The $12$ files ($9$ proof modules, \texttt{Challenge.lean},
\texttt{Proof.lean} and \texttt{CheckSemigroup.lean}) were compiled module by
module in two new empty directories on the author's Mac, with the pinned
toolchain, on 18 November 2026 (23:55--23:59~UTC), and every module was
replayed with \texttt{leanchecker}; each directory also ran one build of the
whole package. All $50$ commands ($25$ per directory) ended with exit
code~$0$, the \texttt{leanchecker} runs printed nothing, and the two
directories produced identical \texttt{.olean} files. The final reviewer (see
the next section) compiled the $12$ files again in a third new directory on
19 November 2026 (01:36--01:41~UTC), with warnings treated as errors, and
replayed every module with \texttt{leanchecker}; all runs ended with exit
code~$0$. Mathlib and its dependencies were taken from a shared prebuilt cache
at the pinned commits and were not rechecked; the \texttt{leanchecker} replays
cover the $12$ project modules. In the repository the files have module names under
\texttt{Research.TaskL.Semigroup} and differ from the checked files only in
their \texttt{import} lines (\texttt{Challenge.lean} is unchanged); these
versions were compiled and replayed with \texttt{leanchecker} again, module by
module, in two new empty directories on the author's Mac (19 November 2026,
02:11--02:17~UTC), with identical \texttt{.olean} files and the same axiom
lists. In the whole-project rebuild of version~1.24.0 of the repository (19--20 November 2026, on the author's Windows computer), these files were compiled and replayed with \texttt{leanchecker} again, together with the rest of the project, each at the first attempt; runs of other modules that failed for lack of memory or file access, without a Lean verdict, were run again and passed (details in the repository). The
formalization has not been checked with comparator or with an independent
kernel.

\par}

\section*{How this paper was produced}
This work was carried out from 18 to 20 November 2026 (UTC) in a workflow
directed by the author, using OpenAI Codex with the model GPT-6.2-Astra
(reasoning level ultra) and Anthropic Claude Code with the model Claude
Opus~5.5. The author asked for a search for substantial open problems, and
Claude Code passed this request to Codex. On 18 November, between 19:22 and
19:50~UTC, Codex screened $50$ problems from the literature, wrote proof
sketches for three of them, had the proofs of this problem read by two other
Codex agents (same model) and checked the statements by exact computation for
all $1249$ numerical semigroups other than $\N$ with Frobenius number at most
$17$; the author chose this problem and one other. From 23:38~UTC, Codex agents
(same model) checked the reading of the sources, searched again for prior
work, wrote the statement file, the written proof and the Lean proof, and ran
the checks. The reading of the sources was checked, and the statement file
reviewed before it was frozen, by one Codex agent that had not taken part in
the search for problems and had not written the statement; this agent then
wrote some of the proof files. After the checks, an independent Claude Code
agent (model Claude Opus~5.5) that had written none of the files gave the
final review. It read the sources itself, compared the statement file with
them, read the written proof, the proof files and the raw logs of the checks,
verified the example of Remark~\ref{rem:sg} by hand and rebuilt the
formalization as described above. Its decision was APPROVE\_FORMAL, with six
minor remarks on the written account, which were taken into account here. According to the
session records, the proof files were written after the freeze by several Codex agents working in
parallel; their compilation logs cover 23:49--23:55~UTC, and the written proof
they follow had been prepared before, in the search for problems. Claude Code (Claude Opus~5.5) checked the sources and
wrote this text from Codex's manuscript and the Lean files. An independent Claude Code agent (model Claude Opus~5.5) that had written none of the files reviewed this paper, and another reviewed the release; their corrections were applied.
Throughout, ``we'' refers to this workflow; the reading and searching described
in the paper were done by the agents. Codex agents searched for prior work on 18 November 2026, and Claude Code searched again on 19 and 20 November 2026; they found no proof of either conjecture beyond multiplicity~$6$ (details in the repository). This reports the coverage of the search, not a guarantee that it was complete. The author is not a
professional mathematician. The author has read this paper and, using a side-by-side table,
compared its Lean statements with Theorems~\ref{thm:interval} and~\ref{thm:even}
and Corollary~\ref{cor:descent}. No human mathematician has reviewed
this paper. The correctness of the results rests on the Lean proofs and on the
agreement of the Lean statements with the conjectures of~\cite{go}; the
written proofs are given for the reader. AI systems are not authors of this
paper; the author takes full responsibility for its content.

\medskip
\noindent Thanks are due to Pedro A. García-Sánchez and Christopher O'Neill
for the questions.

\begin{thebibliography}{9}
\bibitem{br} V.~Blanco and J.~C. Rosales, The tree of irreducible numerical
semigroups with fixed Frobenius number, \emph{Forum Math.} 25 (2013),
1249--1261; arXiv:1105.2147v1 (2011). Proposition numbers refer to the arXiv
version.
\bibitem{bsf} T.~Bogart and S.~A. Seyed Fakhari, Unboundedness of irreducible
decompositions of numerical semigroups, \emph{Comm. Algebra} 53 (2025),
1361--1364; arXiv:2405.11080.
\bibitem{repo} H.~Dong, \emph{ai4math-results}, repository,
\url{https://github.com/Horace-Maxwell/ai4math-results}, version~\repoversion{}
(2026); archived at Zenodo under the concept DOI
\url{https://doi.org/10.5281/zenodo.22970254}.
\bibitem{g} P.~A. García-Sánchez, Factorizations into irreducible numerical
semigroups, \emph{Commun. Korean Math. Soc.} 40 (2025), no.~3, 587--592,
\url{https://doi.org/10.4134/CKMS.c240257}; arXiv:2410.13729.
\bibitem{go} P.~A. García-Sánchez and C.~O'Neill, Lengths of irreducible
decompositions of numerical semigroups, arXiv:2602.00404v1, 2026.
\bibitem{mathlib}
The mathlib Community, The Lean mathematical library, in \emph{Proc.\ 9th ACM
SIGPLAN Int.\ Conf.\ Certified Programs and Proofs} (CPP '20), ACM, 2020,
\url{https://doi.org/10.1145/3372885.3373824}.
\bibitem{rb} J.~C. Rosales and M.~B. Branco, Decomposition of a numerical
semigroup as an intersection of irreducible numerical semigroups,
\emph{Bull. Belg. Math. Soc. Simon Stevin} 9 (2002), 373--381.
\bibitem{rg} J.~C. Rosales and P.~A. García-Sánchez, \emph{Numerical
Semigroups}, Developments in Mathematics~20, Springer, New York, 2009,
\url{https://doi.org/10.1007/978-1-4419-0160-6}.
\end{thebibliography}
\end{document}
